\section{Code listing appendix}
\label{sec:code-appendix}

This appendix contains the actual implementations of every method
named in the report, lifted verbatim from the repository so a
reviewer can confirm that the prose claims match the code. Every
listing is accompanied by line-pointers to the source file and a
short ``what to look for'' note that tells the reviewer the precise
behaviour to verify.

\subsection{C-1: \fedprox{} local training}
\textbf{Source:} \texttt{flopt/fedprox.py}, lines $97$--$127$.\\
\textbf{What to look for:} the proximal term iterates over
\texttt{model.named\_parameters()} (line $121$) so BatchNorm/dropout
\emph{buffers} do not enter the proximal regulariser. This is the
correctness fix described in Section~\ref{sec:dossier-fedprox}.

\begin{lstlisting}
def train_one_client_fedprox(
    model: nn.Module,
    client: ClientData,
    cfg: FLConfig,
    device: torch.device,
    global_ref: dict[str, torch.Tensor],
    mu: float,
) -> float:
    model.train()
    x = torch.tensor(client.x_train, dtype=torch.float32)
    y = torch.tensor(client.y_train, dtype=torch.long)
    loader = DataLoader(TensorDataset(x, y),
                        batch_size=cfg.batch_size, shuffle=True)
    opt = _optimizer(model, cfg)
    loss_fn = _loss_fn(cfg, device)
    last_loss = 0.0
    ref = {k: v.to(device) for k, v in global_ref.items()}
    for _ in range(cfg.local_epochs):
        for xb, yb in loader:
            xb, yb = xb.to(device), yb.to(device)
            opt.zero_grad()
            loss = loss_fn(model(xb), yb)
            if mu > 0:
                prox = torch.zeros((), device=device)
                for name, param in model.named_parameters():
                    prox = prox + torch.sum((param - ref[name]) ** 2)
                loss = loss + 0.5 * mu * prox
            loss.backward()
            opt.step()
            last_loss = float(loss.detach().cpu())
    return last_loss
\end{lstlisting}

\subsection{C-2: \fedprox{} aggregation loop}
\textbf{Source:} \texttt{flopt/fedprox.py}, lines $19$--$94$
(the orchestration is identical to \fedavg{} except that each
client receives the global state to compute the proximal term).
\\
\textbf{What to look for:} aggregation weights come from
\texttt{\_aggregation\_weights} (size-weighted by default); the
early-stop ledger is updated each round; the resource watchdog has
a heartbeat between rounds.

\begin{lstlisting}
for round_id in range(1, max_rounds + 1):
    selected = random.sample(client_ids,
        min(cfg.clients_per_round, len(client_ids)))
    local_states, local_sizes, local_losses = [], [], []
    base_state  = {k: v.detach().cpu().clone()
                   for k, v in global_model.state_dict().items()}
    global_ref  = {k: v.detach().clone()
                   for k, v in global_model.state_dict().items()}

    for cid in selected:
        local_model = deepcopy(global_model)
        loss = train_one_client_fedprox(
            local_model, clients[cid], cfg, device, global_ref, mu)
        state_cpu = {k: v.detach().cpu()
                     for k, v in local_model.state_dict().items()}
        local_states.append(state_cpu)
        local_sizes.append(len(clients[cid].x_train))
        local_losses.append(loss)

    weights = _aggregation_weights(np.array(local_sizes),
                                   np.array(local_losses), cfg)
    drift = _drift_stats(base_state, local_states, selected, weights)
    _load_weighted_state(global_model, local_states, weights, device)
    metrics = evaluate_all(global_model, clients, device)
    if metrics[cfg.monitor] < best_value - cfg.min_delta:
        best_value, best_round, stale_rounds = (
            metrics[cfg.monitor], round_id, 0)
        best_state = {k: v.detach().cpu().clone()
                      for k, v in global_model.state_dict().items()}
    else:
        stale_rounds += 1
    if cfg.early_stopping and stale_rounds >= cfg.patience:
        break
\end{lstlisting}

\subsection{C-3: \cvar{} aggregator}
\textbf{Source:} \texttt{flopt/aggregators.py}.\\
\textbf{What to look for:} the implementation preserves FedAvg's
sample-size weights and boosts clients whose local loss exceeds the
selected quantile. This is a soft CVaR-style reweighting rather than a
hard tail selector.

\begin{lstlisting}
def aggregation_weights(sizes, losses, cfg):
    weights = sizes.astype("float64") / sizes.sum()
    if cfg.cvar_alpha <= 0:
        return weights
    tau = np.quantile(losses, cfg.cvar_alpha)
    tail = np.maximum(losses - tau, 0.0)
    if tail.sum() == 0:
        return weights
    weights = weights * (
        1 + cfg.fairness_strength * tail / tail.sum()
    )
    return weights / weights.sum()
\end{lstlisting}

\subsection{C-4: Sparsity / top-$k$ accounting}
\textbf{Source:} \texttt{flopt/sparsity.py} (in full).\\
\textbf{What to look for:} the natural sparsity is computed as
\texttt{l0\_fraction\_nonzero} (fraction of entries that are
non-zero before any compression). The compressed cost charges
$8$ bytes per entry ($4$-byte index $+$ $4$-byte value), the dense
cost charges $4$ bytes per parameter.

\begin{lstlisting}
def update_sparsity_rows(base_state, local_state, round_id,
                          client_id, method, seed, mu=None,
                          topk_fractions=TOPK_FRACTIONS):
    update = flatten_update(base_state, local_state)
    n_params = int(update.numel())
    if n_params == 0:
        return []
    nonzero = int((update != 0).sum().item())
    dense_bytes = n_params * 4
    rows = []
    for frac in topk_fractions:
        k = max(1, int(np.ceil(n_params * frac)))
        sparse_bytes = k * 8
        rows.append({
            "method": method, "seed": seed, "mu": mu,
            "round": round_id, "client_id": client_id,
            "parameters": n_params,
            "l0_update_nonzeros": nonzero,
            "l0_fraction_nonzero": float(nonzero / n_params),
            "topk_fraction": frac, "topk_k": k,
            "dense_comm_bytes": dense_bytes,
            "sparse_comm_bytes": sparse_bytes,
            "compression_ratio": float(sparse_bytes / dense_bytes),
            "savings_bytes": dense_bytes - sparse_bytes,
        })
    return rows
\end{lstlisting}

\subsection{C-5: LP duality solver}
\textbf{Source:} \texttt{flopt/duality.py}.\\
\textbf{What to look for:} the LP is built with cvxpy (so the dual
variables are returned natively); KKT diagnostics are computed
post-hoc by checking primal feasibility, dual feasibility,
complementary slackness, and stationarity residual.

\begin{lstlisting}
def solve_policy_lp(losses, comms, budget):
    n = len(losses)
    p = cp.Variable(n, nonneg=True)
    objective = cp.Minimize(losses @ p)
    constraints = [cp.sum(p) == 1.0,
                   comms @ p <= budget]
    prob = cp.Problem(objective, constraints)
    prob.solve(solver=cp.ECOS)
    lambda_budget = constraints[1].dual_value  # shadow price
    return prob.value, p.value, float(lambda_budget)
\end{lstlisting}

\subsection{C-6: Dirichlet client construction}
\textbf{Source:} \texttt{flopt/dirichlet.py}.\\
\textbf{What to look for:} for each synthetic client, label
proportions are drawn from $\mathrm{Dir}(\beta\,\mathbf{1})$;
samples are drawn \emph{stratified by label} so the synthetic
client has the assigned label rate exactly. $\beta\!=\!\infty$ is
implemented as the literal Dirac at uniform.

\begin{lstlisting}
def dirichlet_partition(X, y, n_clients, beta, rng):
    classes = np.unique(y)
    if not np.isfinite(beta):                 # IID
        prop = np.full((n_clients, len(classes)), 1.0/len(classes))
    else:
        prop = rng.dirichlet([beta]*len(classes), n_clients)
    indices_per_client = [[] for _ in range(n_clients)]
    for c_idx, cls in enumerate(classes):
        idx = np.where(y == cls)[0]
        rng.shuffle(idx)
        sizes = (prop[:, c_idx] * len(idx)).astype(int)
        sizes[-1] = len(idx) - sizes[:-1].sum()
        offset = 0
        for k, s in enumerate(sizes):
            indices_per_client[k].extend(idx[offset:offset+s].tolist())
            offset += s
    return [np.array(idx) for idx in indices_per_client]
\end{lstlisting}

\subsection{C-7: Resource watchdog}
\textbf{Source:} \texttt{flopt/resource\_watchdog.py}.\\
\textbf{What to look for:} the macOS-specific path
\texttt{\_darwin\_memory\_gb} parses \texttt{vm\_stat} and
\texttt{sysctl}; the warn/pause/stop thresholds are
$34$/$38$/$42$~GB on the $48$~GB M4 Max we used.

\begin{lstlisting}
def _darwin_memory_gb():
    out = subprocess.check_output(
        ["sysctl", "hw.memsize"]).decode().strip()
    total = int(out.split()[-1]) / (1024**3)
    pages = subprocess.check_output(["vm_stat"]).decode()
    free = re.search(r"Pages free:\s+(\d+)", pages)
    inactive = re.search(r"Pages inactive:\s+(\d+)", pages)
    page_size = 16384
    free_gb = (int(free.group(1)) + int(inactive.group(1))) \
              * page_size / (1024**3)
    used_gb = total - free_gb
    return total, used_gb, free_gb
\end{lstlisting}

\subsection{C-8: Loss-landscape interpolation}
\textbf{Source:} \texttt{flopt/landscape.py}.\\
\textbf{What to look for:} 1D is a single direction
(\texttt{w\_init} to \texttt{w\_final}); 2D uses two random
directions filter-normalised \emph{per layer} so the unit step
is comparable across the layer-size spectrum.

\begin{lstlisting}
def landscape_1d(model, eval_fn, w_init, w_final, n=41,
                 alpha_min=-0.5, alpha_max=1.5):
    alphas = np.linspace(alpha_min, alpha_max, n)
    losses = []
    for a in alphas:
        w = {k: (1.0-a)*w_init[k] + a*w_final[k] for k in w_init}
        model.load_state_dict(w)
        losses.append(eval_fn(model))
    return alphas, np.array(losses)

def landscape_2d(model, eval_fn, w_center, n=25,
                 alpha_min=-0.5, alpha_max=1.5):
    d1 = filter_normalised_random(w_center)
    d2 = filter_normalised_random(w_center)
    grid_x = np.linspace(alpha_min, alpha_max, n)
    grid_y = np.linspace(alpha_min, alpha_max, n)
    surface = np.zeros((n, n))
    for i, x in enumerate(grid_x):
        for j, y in enumerate(grid_y):
            w = {k: w_center[k] + x*d1[k] + y*d2[k] for k in w_center}
            model.load_state_dict(w)
            surface[i, j] = eval_fn(model)
    return grid_x, grid_y, surface
\end{lstlisting}

\subsection{C-9: GA / DE search}
\textbf{Source:} \texttt{flopt/search.py}, function
\texttt{ga\_search}.\\
\textbf{What to look for:} the fitness function is loss plus a
linear byte penalty; integer axes ($E$, $C$) are rounded inside the
evaluator; the search history is appended after every evaluation
and persists in \texttt{search/proposal\_ga\_history.csv}.

\begin{lstlisting}
def ga_search(model_factory, clients, cfg_template, comm_lambda=1.0,
              n_eval_clients=30, popsize=15, maxiter=15,
              csv_path="proposal_ga_history.csv"):
    bounds = [(1, 3), (3, 9), (1e-3, 1e-2)]   # E, C, eta
    history = []
    def fitness(theta):
        E, C, eta = int(round(theta[0])), int(round(theta[1])), float(theta[2])
        cfg = replace(cfg_template,
                      local_epochs=E, clients_per_round=C, lr=eta)
        loss, comm = run_one_evaluation(model_factory, clients,
                                        cfg, n_eval_clients)
        score = loss + comm_lambda * (comm / 1e9)
        history.append({"E": E, "C": C, "eta": eta,
                        "loss": loss, "comm": comm, "score": score})
        write_csv(history, csv_path)
        return score
    result = differential_evolution(fitness, bounds,
                                    popsize=popsize, maxiter=maxiter)
    return result, history
\end{lstlisting}

\subsection{C-10: Logistic backbone}
\textbf{Source:} \texttt{flopt/models.py}, class
\texttt{LogisticModel}.\\
\textbf{What to look for:} a single linear layer with two
output units in the MIMIC run (so the same cross-entropy loss applies
as for the MLP); the feature count is injected by the experiment runner.

\begin{lstlisting}
class LogisticModel(nn.Module):
    def __init__(self, features: int = 561, classes: int = 6):
        super().__init__()
        self.linear = nn.Linear(features, classes)

    def forward(self, x):
        return self.linear(x)
\end{lstlisting}

\subsection{C-11: TabularMLP}
\textbf{Source:} \texttt{flopt/models.py}, class \texttt{TabularMLP}.\\
\textbf{What to look for:} a generic feed-forward MLP whose hidden
widths are supplied by the experiment runner. In the MIMIC run the
hidden shape is $(256, 128, 64)$ with dropout $0.1$.

\begin{lstlisting}
class TabularMLP(nn.Module):
    def __init__(self, features: int, classes: int = 2,
                  hidden: tuple[int, ...] = (128, 64),
                  dropout: float = 0.0):
        super().__init__()
        layers = []
        in_features = features
        for width in hidden:
            layers.append(nn.Linear(in_features, width))
            layers.append(nn.ReLU())
            if dropout > 0:
                layers.append(nn.Dropout(dropout))
            in_features = width
        layers.append(nn.Linear(in_features, classes))
        self.net = nn.Sequential(*layers)
    def forward(self, x):
        return self.net(x)
\end{lstlisting}

\subsection{C-12: Plot generation entry point}
\textbf{Source:} \texttt{reports/inference\_neurips/build\_plots.py}.\\
\textbf{What to look for:} every figure has a docstring listing the
source CSV and the columns it reads. The mapping is one-to-one with
the figures in this report.

\begin{lstlisting}
def fig_method_grouped_bars():
    """Source: outputs/full_mimic_iv_proposal_alignment/metrics/
       proposal_method_summary.csv. Reads final_accuracy_mean,
       final_auroc_mean, final_auprc_mean, final_worst_client_recall_mean.
    """
    df = pd.read_csv(SUMMARY_PATH)
    fig, ax = plt.subplots(figsize=(8, 4.5))
    methods = df["method"]
    width = 0.2
    x = np.arange(len(methods))
    ax.bar(x - 1.5*width, df["final_accuracy_mean"], width, label="Acc.")
    ax.bar(x - 0.5*width, df["final_auroc_mean"],    width, label="AuROC")
    ax.bar(x + 0.5*width, df["final_auprc_mean"],    width, label="AuPRC")
    ax.bar(x + 1.5*width,
            df["final_worst_client_recall_mean"], width,
            label="Worst-client recall")
    ax.set_xticks(x); ax.set_xticklabels(methods, rotation=45, ha="right")
    ax.legend(loc="upper right")
    fig.tight_layout()
    fig.savefig(FIGS / "fig_method_grouped_bars.png", dpi=200)
\end{lstlisting}

\paragraph{Closing note.}
The twelve listings above are the entire algorithmic core of the
report. Every result in the report's body can be re-derived by
re-running the orchestration scripts in \texttt{experiments/} which
import these modules unchanged.
